\section{Le calcul mécanisé~: l'outil isolé (1820--1914)}\label{sec:calculating}
% =============================================================================

Pendant que le télégraphe s'électrifie et se déploie le long des voies ferrées, une autre fonction informationnelle croît dans l'ombre~--- mais elle ne s'électrifie pas, ne sort pas de l'organisation qui la commande, et ne crée aucun marché. Les mêmes décennies qui voient Morse breveter son télégraphe (1837), Cooke installer sa ligne sur le \emph{Great Western Railway} (1839) et le câble transatlantique unifier les cours du coton (1866) voient aussi Thomas de Colmar breveter l'Arithmomètre (1820), Babbage concevoir l'\emph{Analytical Engine} (1834) et le \emph{Census Bureau} de Washington crouler sous le dépouillement du recensement de~1880. Les ateliers de mécanique de précision sont parfois les mêmes~: Bréguet, à Paris, fabrique des appareils télégraphiques pour l'administration des lignes et des instruments scientifiques pour les observatoires\footnote{La maison Bréguet, fondée par Abraham-Louis Bréguet en~1775 comme atelier d'horlogerie, fabrique sous son petit-fils Louis-François-Clément les appareils télégraphiques commandés par l'administration française à partir des années~1840 \parencite{bertho_histoire_1984}. Le même atelier produit des chronomètres de marine, des instruments de mesure électrique et des pièces d'horlogerie de haute précision. La mécanique fine du télégraphe et celle de la calculatrice sortent du même monde artisanal.}. Mais les trajectoires divergent~: le télégraphe s'électrifie dès~1837 parce que sa fonction~--- transmettre un signal à distance~--- exige un vecteur que seule l'électricité peut fournir~; le calcul reste mécanique, puis électromécanique, jusqu'à l'arrivée de l'électronique en~1946~: sa fonction~--- combiner des nombres~--- peut être accomplie par des engrenages, des leviers et des doigts humains, puis par des relais électromagnétiques, sans appeler le passage à un substrat radicalement différent. Cette asymétrie~--- la transmission s'électrifie un siècle avant le calcul~--- n'est pas un accident de calendrier~; elle tient à ce que la transmission a une demande externe (le commerce, l'armée, les chemins de fer) qui finance le réseau, tandis que le calcul a une demande interne (la comptabilité, l'actuariat, l'astronomie) qui ne justifie que des outils isolés.

Le calcul est, au dix-neuvième siècle, l'un des travaux les plus répandus et les moins visibles de l'économie industrielle. Les compagnies d'assurance, les administrations fiscales, les chemins de fer emploient des commis par centaines dont le métier est d'additionner des colonnes, de vérifier des sommes, de recopier des résultats d'un registre à l'autre~: un travail répétitif, qualifié, et entièrement interne, dont les chiffres ne sortent pas de l'organisation qui les a commandés. Au moment où les premières machines apparaissent pour l'accélérer, cette activité n'a pas de nom\footnote{Ce n'est qu'en~1962 que \textcite{machlup_production_1962} constitue le traitement de l'information en objet d'analyse économique, sous le nom de \og production de la connaissance\fg{}. Le calcul de bureau du dix-neuvième siècle n'apparaît dans aucune des catégories comptables ou statistiques de l'époque.}. La tension, pourtant, est ancienne. En~1757, le mathématicien Alexis-Claude Clairaut\footnote{Alexis-Claude Clairaut (1713--1765), mathématicien français, membre de l'Académie des sciences à dix-huit ans. Il avait déjà résolu le problème de la figure de la Terre (1743) et clarifié la théorie de la Lune. La prédiction du retour de Halley était pour lui une confirmation expérimentale de la gravitation universelle.} voulut prédire le retour de la comète de Halley en résolvant numériquement le problème à trois corps~: le calcul des perturbations gravitationnelles de Jupiter et de Saturne sur l'orbite de la comète, un test décisif de la mécanique newtonienne que l'astronomie attendait depuis un demi-siècle. Le travail dépassait les forces d'un seul homme~; Clairaut recruta l'astronome Jérôme Lalande\footnote{Jérôme Lalande (1732--1807), astronome français, professeur au Collège de France et directeur de l'Observatoire de Paris. Il sera l'un des astronomes les plus célèbres de son temps~; en~1757, à vingt-cinq ans, il est encore un jeune collaborateur de Clairaut.} et la mathématicienne Nicole-Reine Lepaute\footnote{Nicole-Reine Lepaute (1723--1788), née Étable, épouse de l'horloger Jean-André Lepaute. Mathématicienne et astronome, elle calcule des éphémérides, des tables du Soleil et contribue à la prédiction d'éclipses. Son rôle dans le calcul de la comète sera effacé par Clairaut lui-même (voir la note suivante).}, et les trois s'installèrent au Palais du Luxembourg pour décomposer le problème en pas élémentaires calculés à la plume. Lalande écrira~: \og \emph{Pendant plus de six mois, nous calculâmes depuis le matin jusqu'au soir, quelquefois même à table}\fg{} \parencite[p.~678]{lalande_bibliographie_1803}\footnote{C'est le premier cas documenté de division du travail de calcul entre des personnes recrutées spécifiquement pour cette tâche \parencite[ch.~1]{grier_when_2005}. Clairaut, en publiant sa \emph{Théorie des comètes} (1760), omit de mentionner Lepaute parmi les calculateurs. Lalande lui-même, dans la \emph{Bibliographie astronomique} qu'il fit paraître quarante ans plus tard, rétablit le rôle de Lepaute et critiqua publiquement l'omission~; la querelle avait brisé son amitié avec Clairaut \parencite{lalande_bibliographie_1803}. L'épisode anticipe un trait qui traversera toute la section~: le travail de calcul est le premier à être effacé, et l'effacement est un acte~--- non un oubli~--- qui rend invisibles, en priorité, les contributrices.}. En novembre~1758, Clairaut annonça le retour de la comète pour le 13~avril de l'année suivante. La comète passa à son périhélie le 13~mars~1759, un mois avant la date prédite~: trente et un jours d'écart sur une orbite de soixante-seize ans, soit une erreur inférieure à cinq pour cent. Le calcul avait fonctionné. Il prouvait que les lois de Newton suffisaient à prédire le mouvement d'un objet perturbé par deux planètes, et que six mois de labeur arithmétique pouvaient produire un résultat que l'observation confirmait. D'Alembert, coéditeur de l'\emph{Encyclopédie} et rival de Clairaut à l'Académie des sciences, s'empara de l'erreur résiduelle pour dénoncer ce qu'il appelait \og l'esprit de calcul\fg{} et juger l'entreprise \og plus laborieuse que profonde\fg{}\footnote{La source primaire est \textcite[t.~1, douzième mémoire]{alembert_opuscules_1761}. La querelle est rapportée par \textcite[ch.~1]{grier_when_2005}. Cette dépréciation du calcul \emph{numérique} (résolution itérative, parfois en équipe) au profit du calcul \emph{analytique} (équations différentielles closes, attribuables à un seul auteur) n'est pas une opinion isolée de d'Alembert~: elle structure une part de la culture mathématique académique du XIX\textsuperscript{e}~siècle, jusqu'à la réhabilitation des méthodes numériques au XX\textsuperscript{e}~--- précisément quand la mécanisation du calcul leur donnera un substrat industriel.}. La formule est révélatrice par ce qu'elle refuse de voir~: Clairaut avait divisé le travail intellectuel entre trois personnes, conçu un protocole de calcul séquentiel, et obtenu un résultat que la nature avait confirmé. Mais pour d'Alembert, le labeur numérique ne pouvait pas constituer une méthode. C'est cette tension entre le calcul comme travail et le calcul comme pensée qui traverse la présente section~: qui calcule, à quel coût, et selon quelle organisation~?


\subsection{La demande~: un travail en expansion}\label{subsec:demande_calcul}

En~1774, à Hanovre, un certain Anton Dies, registraire du fonds de pension pour veuves du Calenberg, écrivit au comité de la Trésorerie pour demander une augmentation. Il tenait les comptes du fonds depuis vingt ans, calculait les intérêts composés année par année, vérifiait les tables actuarielles, gérait la correspondance avec les souscripteurs. Lorsqu'une crise des pensions éclata en~1780, le volume de calcul explosa~: le fonds comptait alors plus de 3\,700~couples et 700~veuves, et la réforme exigeait des projections constamment révisées. Dies souffrait d'une paralysie du bras droit. Son successeur Eisendecker estimait sa charge à trois fois celle de son prédécesseur. Un troisième commis, Rath, secondé pour absorber le surplus, sombra dans ce que sa femme qualifia de \og \emph{Gemüthskrankheit}\fg{}, un désordre mental. L'administration lui accorda une gratification de onze thalers\footnote{\textcite{rosenhaft_hands_2003} a reconstitué ce cas à partir des archives de la Landschaft de Hanovre (HStAHann, Dep.~7B et Hann.~93). Les détails biographiques et les montants proviennent de la correspondance administrative entre le comité de la Trésorerie et la Couronne, 1766--1785.}. \dimmark{D5}{?}%
Le cas du Calenberg est modeste, mais il rend visible ce que les comptes d'exploitation du siècle suivant ne montreront jamais~: le calcul est un travail qui use le corps et l'esprit, et dont le volume croît avec la complexité de l'organisation qu'il sert. Reformulée dans le vocabulaire marshallien~--- que les administrateurs hanovriens ne mobilisaient évidemment pas~---, la structure du cas peut être résumée ainsi~: la demande de calcul est dérivée (on calcule pour assurer, pour dénombrer, pour acheminer, jamais pour calculer), peu sensible au prix de la machine~--- un commis coûte le même salaire que la machine existe ou non, et l'investissement dans une machine ne réduit pas le calcul que l'organisation doit produire~---, et croissante avec la complexité de l'organisation servie~: plus le nombre de souscripteurs, de polices, de wagons, de courses croît, plus le volume de calcul à exécuter croît. Ces propriétés sont qualitatives, reconstituées à partir d'un cas individuel~; elles ne prétendent pas à l'estimation d'élasticités, qui supposerait des séries de prix et de quantités dont les archives administratives hanovriennes ne disposaient pas. \dimmark{0-TER}{?}%
Ce que les sources documentent, c'est l'usure du corps et de l'esprit~: le bras paralysé de Dies, la \og \emph{Gemüthskrankheit}\fg{} de Rath, l'épuisement d'Eisendecker. Aucun poste budgétaire n'enregistre ces coûts, et cette absence n'est pas seulement empirique~: elle est constitutive. La catégorie même d'\emph{énergie} comme grandeur unifiée n'existait pas en~1780~--- elle ne se stabilisera qu'avec l'unification thermodynamique du milieu du XIX\textsuperscript{e}~siècle (Mayer, Joule, Helmholtz)~---, et son importation dans la pensée économique sous la forme d'un travail humain converti en équivalent énergétique restera marginale, depuis les tentatives de Sergei Podolinsky (1880) jusqu'à leur reprise tardive par Nicholas Georgescu-Roegen un siècle plus tard\footnote{Sur la généalogie de l'\emph{énergie} comme catégorie économique~--- de Podolinsky, Sacher, Bogdanov, Ostwald et Neurath jusqu'aux fondateurs de l'économie écologique au XX\textsuperscript{e}~siècle~---, on renverra à Martinez-Alier, \emph{Ecological Economics. Energy, Environment and Society} (Oxford, Blackwell, 1987), et à la reconstruction récente de \textcite{vianna_franco_history_2023}. La présente section mobilise quelques jalons de cette histoire sans en faire un objet pour lui-même.}. Que le calcul du Calenberg n'apparaisse pas comme \og dépense énergétique\fg{} ne tient donc pas à un oubli des comptables hanovriens~: le concept qui aurait permis de l'enregistrer comme tel n'était pas encore disponible. À l'invisibilité d'ordre de grandeur~--- le coût métabolique d'un commis est dérisoire devant le système qu'il sert~--- s'ajoute donc une invisibilité catégorielle plus profonde~: un coût qui n'a pas de nom n'a pas de poste.

\dimmark{D4}{?}%
Cette invisibilité n'est pas accidentelle. \textcite{downey_virtual_2001} a montré que les discussions sur la \og société de l'information\fg{}, contemporaines comme rétrospectives, effacent systématiquement le travail humain dans les réseaux d'information. Le travail de calcul en bureau n'a pas laissé de traces comparables aux registres commerciaux ou aux brevets d'invention~; les budgets des entreprises du dix-neuvième siècle, quand ils sont conservés, ne distinguent pas la ligne \og calcul\fg{} de la masse salariale administrative. \textcite{rosenhaft_hands_2003} suggère que cette invisibilité tient en partie au caractère \og \emph{low-tech}\fg{} du travail~: à la différence des télégraphistes ou des opératrices téléphoniques, dont le travail était visiblement lié à un appareil, les commis qui additionnaient des colonnes et recopiaient des totaux ne semblaient pas exercer une activité qui fût en prise avec la technique.

On peut néanmoins en mesurer la croissance. Le nombre de commis aux États-Unis passe de un à deux pour cent de la population active en~1870 à plus de dix pour cent en~1940 \parencite{cortada_before_1993}\footnote{La catégorie \og clerk\fg{} dans les recensements américains a connu plusieurs reclassifications entre~1880 et~1930. \textcite{cortada_before_1993} construit ses séries en harmonisant autant que possible les nomenclatures successives, mais les comparaisons longues sur sept décennies restent sensibles à ces reclassifications et sont à manier avec précaution.}. \textcite{rotella_transformation_1981} a établi, à partir des données du Census Bureau, la correspondance temporelle précise entre l'adoption des machines de bureau et le changement de composition par genre de l'emploi clérical entre~1870 et~1930. En Angleterre, \textcite{boot_salaries_1991} a montré que le nombre de clercs de la Bank of Scotland plus que tripla entre~1730 et~1880, alors que le nombre de guichetiers ne fit que doubler. \dimmark{D7}{+}%
L'écart de croissance entre le \og \emph{front office}\fg{} et le \og \emph{back office}\fg{} suggère un effet multiplicateur interne~: chaque augmentation du volume d'affaires engendrait un accroissement plus que proportionnel du travail de tenue des registres, de vérification et de correspondance. C'est dans cet écart que les premières machines vont s'insérer. Mais le fait remarquable est la lenteur de cette insertion~: la première machine capable d'effectuer les quatre opérations arithmétiques est brevetée en~1820, et il faudra attendre un demi-siècle pour que la demande soit suffisamment concentrée pour en absorber la production.

\emergence{D5/D6/0-TER~: le calcul mécanisé diffère structurellement du télégraphe sur trois dimensions. Premièrement, la demande de calcul est dérivée et peu sensible au prix de la machine (D5)~: là où la baisse du prix du télégramme augmente le trafic, la baisse du prix de la calculatrice ne crée pas de demande nouvelle de calcul~; c'est la complexité de l'organisation servie qui détermine le volume. Deuxièmement, les effets du calcul restent internes à l'organisation~: la calculatrice ne crée pas de marchés, ne synchronise pas de prix, ne rend pas possible des transactions qui n'auraient pas eu lieu sans elle (D6 négatif, par contraste avec le résultat de Steinwender pour le télégraphe). Troisièmement, l'énergie du calcul est métabolique~: elle passe par le corps du commis, et à ce titre elle est doublement invisible~: invisible dans les comptes parce qu'elle ne constitue pas un poste identifié, et invisible dans les ordres de grandeur parce qu'elle est dérisoire comparée au système qu'elle sert (0-TER). \questionch{2}{D5/D6~: l'effet des techniques de l'information sur les marchés dépend-il de la fonction exercée (transmission, stockage ou calcul)~? Si oui, la codifiabilité de Juhász n'est pas le seul filtre~: la nature même de la fonction ICT en est un autre.}}


\subsection{De l'Arithmomètre aux calculatrices industrielles}\label{subsec:arithmometre}

En~1820, Charles Xavier Thomas de Colmar\footnote{Thomas de Colmar (1785--1870), né Charles Xavier Thomas, directeur de la compagnie d'assurance Le~Phénix à Paris. Son brevet d'invention (n\textsuperscript{o}~1420, 18~novembre 1820, archives INPI) décrit un \og arithmomètre ou machine à calculer\fg{}. L'appareil est décrit dans le catalogue officiel de la \emph{Great Exhibition} de Londres en~1851, où il reçut une médaille, et dans la notice technique de \textcite{sebert_machines_1879}.}, directeur de la compagnie d'assurance Le~Phénix à Paris, déposa un brevet pour un appareil qu'il appelait \og arithmomètre\fg{}~--- une machine capable d'effectuer les quatre opérations arithmétiques par un mécanisme de tambours cannelés\footnote{Le tambour cannelé (\emph{Staffelwalze} ou \emph{Leibniz wheel}) est un cylindre portant neuf dents de longueur croissante, conçu par Leibniz pour son \emph{Stepped Reckoner} (1694). Le principe permet de réaliser la multiplication par rotation différentielle. Leibniz en donna une description dans les \emph{Miscellanea Berolinensia} \parencite{leibniz_machina_1710}. Le mécanisme resta sans descendance industrielle pendant plus d'un siècle, jusqu'à ce que Thomas de Colmar le reprenne dans un boîtier de dimensions compatibles avec un bureau.} dérivé du \emph{Stepped Reckoner} de Leibniz. Que la première machine à calculer produite en série fût l'\oe{}uvre d'un assureur n'est pas un hasard~: le calcul actuariel~--- tables de mortalité, primes, réserves~--- était l'un des rares domaines où le volume de calcul justifiait l'investissement dans un appareil coûteux et lent. L'Arithmomètre fonctionnait~; il fonctionnait assez bien pour que des observatoires et quelques administrations l'achètent à leur tour, et assez longtemps pour rester en service jusqu'à la Première Guerre mondiale \parencite{cortada_before_1993, sebert_machines_1879}\footnote{La machine n'était pas équipée d'imprimante~; pour pallier cette absence et garantir la fiabilité de la transcription des résultats, \textcite{sebert_machines_1879} recommandait dans son rapport à la Société d'encouragement pour l'industrie nationale d'avoir recours à \og deux employés, un qui tourne la manivelle et l'autre qui transcrit\fg{}. La calculatrice ne supprimait pas le travail de calcul~: elle imposait une division du travail à deux personnes autour d'elle, ce qui constitue l'un des premiers cas documentés de réorganisation du travail induite par une machine de bureau~--- pattern que l'on retrouvera, amplifié, autour des tabulatrices Hollerith au \S\,\ref{sec:mechanography}.}. Mais la diffusion fut d'une lenteur remarquable. \textcite{cortada_before_1993} estime que Thomas vendit \og \emph{possibly several hundred}\fg{} machines sur trois à quatre décennies~; \textcite{nordhaus_two_2007} retient un stock cumulé d'environ cinq cents machines avant~1865. La capacité technique précédait la demande~: la machine existait, mais le marché qui aurait pu l'absorber~--- des organisations dont les besoins de calcul fussent suffisamment volumineux et répétitifs pour justifier l'investissement~--- ne se constituerait qu'avec la croissance des compagnies d'assurance et des administrations statistiques dans la seconde moitié du siècle.

En~1834, Charles Babbage conçut l'\emph{Analytical Engine}, une machine dont l'architecture se distinguait radicalement de l'Arithmomètre\footnote{La machine à différences elle-même, sur laquelle Babbage avait travaillé depuis les années~1820, n'était pas conceptuellement nouvelle~: la même proposition avait été faite par un officier hessois, J.~H.~Müller, en~1786, sans que Babbage en eût connaissance \parencite[p.~42]{augarten_bit_1984}. L'idée d'une mécanisation du calcul polynomial par différences finies était disponible dans la culture mathématique européenne depuis près d'un demi-siècle~; ce qui manquait n'était pas le concept mais les conditions économiques de sa réalisation.}. L'appareil séparait deux organes~: le \emph{store}, un ensemble de colonnes de roues capables de retenir cinquante nombres de cinquante chiffres, et le \emph{mill}, une unité de calcul capable de les combiner selon quatre opérations. Les instructions devaient être codées sur des cartes perforées empruntées au métier Jacquard~; le même mécanisme servait à sélectionner les nombres stockés et à commander la séquence des opérations. L'\emph{Analytical Engine} ne fut jamais construit~; mais sa conception posait, pour la première fois dans un artefact technique, la distinction entre le stockage et le traitement de l'information comme deux fonctions architecturalement séparées\footnote{Ada Lovelace, dans ses notes de~1843 sur le mémoire de Menabrea, en décrivit les possibilités avec une précision qui en fait le premier document de programmation de l'histoire. \textcite{schaffer_babbage_1994} souligne un paradoxe~: attribuer l'intelligence à la machine occultait le travail humain qui restait nécessaire pour la concevoir, la programmer et en vérifier les résultats.}.

La machine à différences, en revanche, fut effectivement construite~--- mais en Suède, par un autre que Babbage. Georg Scheutz, juriste et éditeur de journal stockholmois, lut un compte rendu des travaux de Babbage dans l'\emph{Edinburgh Review} de~1834 et entreprit, avec son fils Edvard et l'appui de l'Académie royale des sciences de Suède, de réaliser l'appareil. Une première machine fonctionnait en~1844~; une version industrielle, capable de produire des stéréotypes d'imprimerie pour la composition automatique des tables, fut achevée en~1855 et exposée à l'Exposition universelle de Paris la même année. L'une de ces machines fut achetée par le \emph{General Register Office} britannique et y fonctionna entre~1859 et~1864 pour calculer les tables actuarielles du \emph{English Life Table}~\parencite[ch.~1]{goldstine_computer_1972}\footnote{La machine Scheutz est aujourd'hui conservée au Smithsonian. Elle constitue, à notre connaissance, le seul cas documenté d'utilisation industrielle d'une machine babbagéenne au XIX\textsuperscript{e}~siècle. \textcite{winston_media_1998} en fait l'illustration de ce qu'il appelle un \og prototype accepté en niche\fg{}~: un dispositif techniquement abouti et économiquement utile dans une demande étroite, mais qui n'entraîne pas la généralisation faute d'une nécessité sociale survenante de plus large portée. Les performances effectives de la machine, soulignent toutefois \textcite[p.~328]{warwick_laboratory_1995}, restèrent faibles, en raison notamment de défauts d'impression~--- ce qui n'invalide pas l'argument de niche réalisée mais en nuance la portée.}. \dimmark{D5}{?}%
Le cas Scheutz importe pour notre argument~: il établit que la machine à différences était techniquement réalisable et économiquement utile dans une niche, et que la non-émergence de l'industrie babbagéenne ne tient ni à un verrou technique ni à une impossibilité économique en soi, mais à l'étroitesse de la niche. Le \emph{General Register Office} avait besoin de tables de mortalité~; il ne lui en fallait pas plusieurs centaines par an. La machine Scheutz desservait une demande qui se satisfaisait d'un exemplaire unique. Babbage, lui, projetait un calculateur universel pour une demande qui n'existait pas encore. La distinction entre les deux projets n'est pas seulement technique~--- la machine à différences calcule des polynômes par différences finies tandis que l'\emph{Analytical Engine} aurait pu, par sa programmabilité altérable, traiter n'importe quel problème calculable\footnote{C'est précisément l'intuition de Babbage et Lovelace que reprendra Turing un siècle plus tard~: la programmabilité altérable~--- la capacité d'une machine à modifier son propre programme selon les résultats intermédiaires~--- distingue le calculateur de l'ordinateur \parencite[ch.~8]{winston_media_1998}. Cette distinction reste \og indistincte\fg{} (\emph{indistinct conception}) dans les carnets tardifs de Babbage, ce qui le maintient en deçà du seuil conceptuel franchi en~1936 par \textcite{turing_computable_1937}.}~--- elle est aussi économique~: une niche ne nourrit pas une industrie, et seule une demande répétée en volume élevé pouvait justifier l'investissement qu'exigeait une production sérielle de machines à calculer.

\dimmark{D8}{?}%
L'échec de Babbage n'est pas seulement technique~; il est économique~--- et le paradoxe est que Babbage lui-même avait formulé, dans son \emph{On the Economy of Machinery and Manufactures} (Londres, Charles Knight, 1832), l'une des premières théories systématiques de la mécanisation comme division du travail\footnote{Le \og principe Babbage\fg{}~--- décomposer une tâche complexe en sous-tâches de qualifications hétérogènes, afin de ne payer chaque qualification qu'au prix qu'elle exige~--- est exposé au chapitre~19 de l'\emph{Economy of Machinery}. Marx en reprendra l'argument dans \emph{Le Capital} (livre~I, ch.~14). Babbage disposait donc du cadre analytique qui aurait pu rendre intelligible son propre projet~--- mais ce cadre supposait précisément le marché qu'il n'avait pas.}. Le \emph{Difference Engine} n\textsuperscript{o}~1, conçu pour calculer des tables polynomiales par la méthode des différences finies, avait reçu environ 17\,470~livres sterling du gouvernement britannique entre~1823 et~1842, une somme comparable au coût d'un navire de guerre de l'époque \parencite{swade_difference_2001}. Le projet fut abandonné après que Babbage et son ingénieur Joseph Clement se brouillèrent sur la propriété des outils, et le gouvernement refusa de financer l'\emph{Analytical Engine}. Le contraste avec le câble transatlantique est instructif~: le câble suivait la séquence investissement, faillite, recapitalisation, succès~; la machine de Babbage suivait la séquence investissement, pas de produit, pas de recapitalisation, abandon. La différence tient à la demande~: en~1840, aucune organisation n'avait besoin de calculer en volume au point de justifier un investissement de cette ampleur~; cinquante ans plus tard, le Census Bureau américain en aurait besoin.

Ce que l'échec de Babbage confirmait, c'est ce que l'Arithmomètre suggérait par son succès modeste~: la machine à calculer ne pouvait se déployer que là où la demande de calcul était déjà organisée en flux régulier, et cette condition ne serait remplie, à grande échelle, que dans le dernier quart du siècle.

Le passage de l'Arithmomètre artisanal à la calculatrice industrielle s'accomplit entre~1875 et~1914 par une convergence technique qui mit fin au monopole de fait du tambour de Leibniz. En~1875, Frank Baldwin aux États-Unis et Willgodt Odhner en Russie développèrent indépendamment un mécanisme à disque à broches rétractables~--- plus compact, plus fiable, moins coûteux à fabriquer que le tambour cannelé. L'ajout du clavier, emprunté à la machine à écrire, transforma la calculatrice en un instrument de bureau ergonomique. En une génération, quatre \emph{dominant designs} s'imposèrent~: le Comptometer de Dorr Felt (1886, à clavier direct, sans manivelle), la machine Burroughs (1886, à clavier avec impression du résultat), les machines Odhner et Brunsviga (à manivelle, dominantes en Europe) \parencite{cortada_before_1993}.


\subsection{La diffusion des calculatrices et la structure du marché}\label{subsec:burroughs}

La croissance des ventes, documentée par les archives de la \emph{Burroughs Adding Machine Company} exploitées par \textcite{cortada_before_1993}, fut spectaculaire. L'\emph{American Arithmometer Company} (future Burroughs) vendit 286~machines en~1895 à un prix moyen de 223~dollars~; 1\,399 en~1900~; 5\,008 en~1904~; 15\,763 en~1909~--- le double de ce que William Burroughs pensait, en~1885, pouvoir vendre dans le monde entier. Les ventes annuelles passèrent de 4,9~millions de dollars en~1905 à 39,9~millions en~1916, soit un facteur huit en onze ans~--- une croissance largement supérieure à celle du produit national brut américain sur la même période. En Europe, Brunsviga vendit plus de vingt mille machines entre~1885 et~1912 et obtint 130~brevets allemands et 300~brevets étrangers~; le Millionaire suisse, capable de multiplication directe par lecture d'une table de multiplication mécaniquement stockée dans la machine, plaça 4\,600~machines entre~1894 et~1935 \parencite{cortada_before_1993}\footnote{La table de multiplication du Millionaire, gravée sur un bloc métallique mobile, anticipe matériellement la fonction de mémoire morte (\emph{ROM}) des calculateurs électroniques~: un savoir codifié, fixé dans la matière de la machine, et lu sans recalcul à chaque opération. Ce dispositif illustre que la fonction de stockage~--- ce que la présente thèse appelle STORE~--- n'est pas née avec l'électronique~; elle est constitutive de toute machine à calculer dès qu'elle dépasse l'opération atomique \parencite[ch.~8]{winston_media_1998}. Le principe des \og barèmes matérialisés\fg{}, c'est-à-dire l'inscription mécanique d'une table de multiplication dans la machine elle-même sous forme de chevilles en saillie sur des plaques, avait été développé indépendamment par Léon Bollée au Mans dès~1889~--- précédant Steiger qui ne breveta le Millionaire qu'en~1892~--- et présenté à la Société d'encouragement pour l'industrie nationale par \textcite{sebert_bollee_1890}.}. \textcite{nordhaus_two_2007}, à partir d'une reconstruction patiente de la productivité du calcul sur deux siècles, estime que le stock cumulé mondial de calculatrices mécaniques passa de cinq cents avant~1865 à environ 130\,000 en~1910 et 900\,000 en~1920, et que le coût par opération arithmétique baissa d'environ trois pour cent par an entre~1850 et~1945.

\dimmark{D2}{+}%
Le marché des calculatrices mécaniques ne convergeait pas vers le monopole. Une centaine de firmes fabriquaient des calculatrices aux États-Unis en~1904~; treize fabricants restaient actifs dans les années~1920 \parencite{cortada_before_1993}. Les barrières à l'entrée étaient faibles~: la mécanique de précision ne requérait pas d'infrastructure de réseau, la vente directe ou par agents ne supposait pas de couverture nationale, et la machine ne créait aucune dépendance au fournisseur une fois achetée. \textcite{leffingwell_office_1926}, dans le \emph{Office Appliance Manual} publié par l'association professionnelle du secteur, recensait des dizaines de modèles concurrents vendus à des prix comparables par des canaux similaires~--- un marché fragmenté, disputé, sans concentration visible.

\dimmark{D6}{--}%
L'effet de ces machines sur l'économie restait confiné à l'intérieur des organisations. La calculatrice n'agissait pas sur les prix entre entreprises, ne créait pas de marchés nouveaux, ne modifiait pas les échanges~: elle agissait sur les coûts de fonctionnement, sur la division interne du travail, sur la productivité du commis. \textcite{mills_white_1951}, dans \emph{White Collar}, observait que les machines de bureau ne furent pas utilisées de manière extensive avant~1910~--- une date que \textcite{cortada_before_1993} corrige à la hausse en montrant que les volumes de vente étaient déjà substantiels dans les années~1890. Mais l'un comme l'autre constatent que les effets de ces machines n'apparaissaient pas sur les marchés~: ils restaient dans les budgets internes des organisations, et les budgets des entreprises du dix-neuvième siècle ne sont, pour l'essentiel, pas conservés.


\subsection{La restructuration du travail de bureau}\label{subsec:calcul_travail}

\dimmark{D4}{*}%
Le déclin du coût unitaire du calcul n'entraîna pas de déclin de l'emploi~: au contraire, le nombre de commis aux États-Unis passa de un à deux pour cent de la population active en~1870 à plus de dix pour cent en~1940 \parencite{cortada_before_1993}. \textcite{reinstaller_market_2001}, dans le seul travail qui construise le mécanisme causal complet, ont avancé que le biais technique américain~--- capital-intensif, \emph{labour-saving}~--- transféra à l'intérieur du bureau la même logique de substitution qui avait transformé l'usine~: la corrélation entre le ratio capital/travail global de l'économie américaine et le ratio capital/travail des employés de bureau atteint~0,921 entre~1889 et~1937\footnote{Une corrélation aussi élevée sur deux séries également tendancielles ne suffit pas, en toute rigueur, à établir un mécanisme causal commun~: elle est compatible avec l'effet d'une tendance partagée. Reinstaller \& Hölzl mobilisent cependant l'argument dans une chaîne plus large (biais technique, baisse du salaire réel, féminisation), où la corrélation joue le rôle d'une régularité statistique parmi d'autres pièces.}. La part des femmes dans l'emploi de bureau passa de 2,4\,\% en~1870 à 48,4\,\% en~1920~; les salaires réels des commis baissèrent de 1\,920 à 1\,385~dollars (en prix de~1929) entre~1899 et~1919, les femmes étant payées 25 à 45\,\% de moins que les hommes \parencite{reinstaller_market_2001}. \textcite{davies_womans_1982} a documenté, à partir d'archives d'entreprise et de données du Census Bureau, la dynamique de ce basculement~: les employeurs embauchaient des femmes non parce qu'elles étaient formées aux machines, mais parce que les machines rendaient possible le recrutement d'une main-d'\oe{}uvre moins payée et la réorganisation du travail de bureau selon un modèle plus proche de la manufacture. \textcite{strom_beyond_1992} a étendu l'analyse en montrant que la féminisation ne se réduisait pas à un effet de prix~: elle s'accompagnait d'une transformation des hiérarchies, des espaces et des rythmes de travail dans le bureau américain entre~1900 et~1930. La calculatrice ne supprimait pas le travail~: elle le restructurait, le féminisait, et en comprimait le coût~--- non pas en éliminant des postes mais en substituant à des commis qualifiés une main-d'\oe{}uvre moins qualifiée, moins payée, assistée par la machine. \textcite{wootton_making_2007} confirme le résultat par un autre angle~: l'introduction des machines de bureau dans la comptabilité britannique ne réduisit pas le nombre de \emph{bookkeepers} mais modifia la nature de leur tâche, la partie mécanique du calcul étant transférée à la machine et la partie cognitive~--- vérification, interprétation, décision~--- restant humaine.

L'énergie, dans cette histoire, était presque inexistante. La seule donnée quantitative que nous ayons trouvée provient des archives du \emph{Nautical Almanac} exploitées par \textcite{daston_calculation_2018}~: en~1929, l'électricité nécessaire au fonctionnement des machines Hollerith et Burroughs représentait 9~livres sterling sur un budget total de 607~livres, soit 1,5\,\% du coût. Cette donnée est isolée~--- un point dans une institution une année~--- et il convient de la prendre comme illustrative plutôt que comme l'établissement d'une régularité~; nous n'avons pas trouvé de séries comparables sur d'autres budgets de l'époque\footnote{La part électrique des budgets administratifs et bancaires des années~1920--1930 reste un point aveugle de l'historiographie comptable. \textcite{cortada_before_1993} et \textcite{wootton_making_2007} ne discutent pas spécifiquement le poste \og énergie\fg{} dans les budgets qu'ils exploitent. Le résultat de \textcite{daston_calculation_2018} ne peut donc pas être généralisé sans réserve, mais il indique un ordre de grandeur cohérent avec ce que l'on attend d'un calcul électromécanique d'avant l'électronique.}. \dimmark{0-TER}{+}%
Sous cette réserve, le couplage entre calcul et énergie présentait deux propriétés qui le rendaient invisible~: d'abord par les ordres de grandeur~--- l'électricité, lorsqu'elle existait, restait un poste négligeable dans les budgets disponibles~---, ensuite par le fait que le calcul ne pilotait pas directement un système physique. La calculatrice produisait des chiffres qui restaient sur le papier~; elle ne contrôlait pas un réseau, ne coordonnait pas des flux matériels, ne consommait pas de matières premières au-delà du papier et de l'encre. L'industrie entière des machines de bureau~--- calculatrices, tabulatrices, machines à écrire, caisses enregistreuses~--- pesait 217,8~millions de dollars en~1929, soit deux dixièmes de pour cent du produit national brut américain \parencite{cortada_before_1993}. La question énergétique du calcul ne se poserait qu'au moment où le calcul changerait de substrat~--- du mécanique à l'électronique~--- et où son volume croîtrait de plusieurs ordres de grandeur.

\emergence{D2~--- Le marché des calculatrices mécaniques ne converge pas vers le monopole~: une centaine de firmes en~1904, treize fabricants encore actifs dans les années~1920. Les barrières à l'entrée sont faibles, le produit ne crée pas de dépendance au fournisseur, et aucun rendement croissant de réseau ne joue. \questionch{2}{D2~--- La convergence vers le monopole est-elle liée aux propriétés du réseau de communication, ou peut-elle procéder d'un autre mécanisme dans le calcul~?}}

\emergence{D4~--- La calculatrice ne remplace pas le travail~: elle le restructure. L'emploi de bureau croît, se féminise, et le salaire réel baisse. \textcite{davies_womans_1982} et \textcite{strom_beyond_1992} documentent la transformation des hiérarchies et des espaces de travail qui accompagne la mécanisation. \questionch{2}{D4~--- Cette restructuration sans suppression est-elle propre à la calculatrice, ou se retrouve-t-elle avec la tabulatrice et l'ordinateur~?}}

Au-delà des résultats par dimension, ce parcours traverse en filigrane une transformation plus large de l'économie matérielle de l'écrit administratif, dont \textcite{gardey_ecrire_2008} a proposé la périodisation et le cadre interprétatif. Entre~1890 et~1940, l'organisation du travail de bureau bascule d'un régime où la sécurité, la conservation et la traçabilité sont les valeurs cardinales~--- celles qui justifiaient le commis qualifié, l'écriture lente, le registre relié~--- vers un régime où la vitesse, la malléabilité et la mobilité de l'information prennent le pas. La confiance que les pratiques anciennes plaçaient dans les humains et dans leur jugement est progressivement déléguée aux artefacts et aux dispositifs qui en organisent le repérage. Cette bascule normative, que Gardey nomme une \og économie morale\fg{} de l'inscription, n'est pas le simple effet de la mécanisation~: elle en est la condition. La calculatrice, dont les effets restaient internes à l'organisation, accompagne ce déplacement sans le saturer~; la mécanographie qui suit, en couplant le stockage industrialisé et le calcul automatisé, viendra l'inscrire dans un dispositif d'une autre nature. C'est l'objet de la section suivante.


% =============================================================================